\section{Discussion}
\label{sec:discussion}

Taken together, the results point to a layered interpretation of Brazilian equity dependencies, in line with the multi-scale perspective advocated by Raddant and Di Matteo~\cite{raddant_dimatteo_2023}. The first layer is a market-mode component. It is reflected in the moderate positive average correlation ($\bar{\rho}\approx0.24$), the increase in rolling average correlation around stress periods and the very large leading eigenvalue ($\lambda_1=21.65$, approximately 37.3\% of the trace of the correlation matrix). We interpret this component as broad market synchronization. The analysis does not identify the economic mechanisms behind that synchronization; its role here is to show why raw networks can obscure weaker localized patterns.

The second layer consists of candidate group and sector structure. Within-sector correlations are higher than between-sector correlations on average, and the RMT decomposition identifies four additional eigenvalues above the random upper edge. Eigenvector loadings and group-mode matrices suggest that commodity, financial, telecommunications and utility-related structures contribute to the non-random part of the spectrum. This reading is descriptive, not causal: we identify structure in the dependency matrix, not the economic mechanisms that generate each edge.

The third layer is noise-compatible residual dependence. Not every correlation is economically meaningful, and not every edge retained by a graph is a stable economic relationship. This is why we combine RMT, heatmaps, dendrograms and network metrics. A result is more credible when it appears across several representations, not only in one visualization.

The network results show that market-mode removal changes structure in a systematic way. Original networks have stronger average edge correlations, while group-mode networks have much lower average edge correlations. However, the group-mode PMFG has higher clustering (0.6653 versus 0.5273). Notably, lower average edge strength does not mean absence of structure. It means that the dominant common component has been removed and localized organization becomes easier to observe \cite{tumminello_2005_tool,tumminello_2010_correlation}. The PMFG is well suited for this analysis because it preserves richer graph structure than the MST, retaining 168 edges with 166 triangles and 55 4-cliques that allow hub-rank shifts, clustering changes and subsector dependency maps to be studied in detail.

Clustering summarizes hierarchical organization, but it does not identify graph-theoretic hubs, bridges or local cycles. The transition from dendrograms to filtered financial networks adds information. Hub reorganization after RMT filtering is meaningful: original hubs such as GGBR4 and BBDC4 are central in networks that include the common component, while group-mode hubs such as GUAR3 occupy central positions after the leading component is removed. This suggests that centrality in financial networks depends on the matrix used to construct the network.

The volatility-modelling results require caution. Machine-learning and simple deep-learning models do not produce uniform differences across horizons and metrics. Instead, models augmented with ex-post network descriptors are competitive at selected horizons, while lagged volatility variables remain dominant. Because the RMT and network descriptors are estimated from the full sample, these results identify a retrospective structural association rather than a real-time forecast gain. The net interpretation is modest: topology summarizes cross-sectional geometry that may complement the temporal dynamics of individual asset variance, but its point-in-time value remains untested.

The broader contribution is an integrated and reproducible framework connecting econophysics filtering, financial networks and volatility modelling for B3 equities. Separating chronological model fitting from ex-post structural enrichment makes the current evidence easier to interpret and defines the next step: reconstruct RMT and network quantities in rolling or expanding windows at every forecast origin.
